Um die Heisenberg-Gleichung der Bewegung herzuleiten, beginnen wir mit der Definition des Heisenberg-Operators:

\[
\hat{O}_H(t) = e^{\frac{i}{\hbar} \hat{H}(t-t_0)} \hat{O} e^{-\frac{i}{\hbar} \hat{H}(t-t_0)}
\]

Wir berechnen die zeitliche Ableitung von \(\hat{O}_H(t)\):

\[
\frac{d \hat{O}_H}{d t} = \frac{d}{d t} \left( e^{\frac{i}{\hbar} \hat{H}(t-t_0)} \hat{O} e^{-\frac{i}{\hbar} \hat{H}(t-t_0)} \right)
\]

Verwenden der Produktregel ergibt:

\[
\frac{d \hat{O}_H}{d t} = \left( \frac{d}{d t} e^{\frac{i}{\hbar} \hat{H}(t-t_0)} \right) \hat{O} e^{-\frac{i}{\hbar} \hat{H}(t-t_0)} + e^{\frac{i}{\hbar} \hat{H}(t-t_0)} \hat{O} \left( \frac{d}{d t} e^{-\frac{i}{\hbar} \hat{H}(t-t_0)} \right)
\]

Die Ableitungen der Exponentialoperatoren sind:

\[
\frac{d}{d t} e^{\frac{i}{\hbar} \hat{H}(t-t_0)} = \frac{i}{\hbar} \hat{H} e^{\frac{i}{\hbar} \hat{H}(t-t_0)}
\]

\[
\frac{d}{d t} e^{-\frac{i}{\hbar} \hat{H}(t-t_0)} = -\frac{i}{\hbar} \hat{H} e^{-\frac{i}{\hbar} \hat{H}(t-t_0)}
\]

Einsetzen dieser Ausdrücke in die Produktregel ergibt:

\[
\frac{d \hat{O}_H}{d t} = \frac{i}{\hbar} \hat{H} e^{\frac{i}{\hbar} \hat{H}(t-t_0)} \hat{O} e^{-\frac{i}{\hbar} \hat{H}(t-t_0)} - e^{\frac{i}{\hbar} \hat{H}(t-t_0)} \hat{O} \frac{i}{\hbar} \hat{H} e^{-\frac{i}{\hbar} \hat{H}(t-t_0)}
\]

Dies kann umgeschrieben werden als:

\[
\frac{d \hat{O}_H}{d t} = \frac{i}{\hbar} \left( \hat{H} \hat{O}_H(t) - \hat{O}_H(t) \hat{H} \right)
\]

Dies ist der Kommutator:

\[
\frac{d \hat{O}_H}{d t} = \frac{i}{\hbar} [\hat{H}, \hat{O}_H(t)]
\]

Multiplizieren mit \(i \hbar\) ergibt die Heisenberg-Gleichung der Bewegung:

\[
i \hbar \frac{d \hat{O}_H}{d t} = [\hat{O}_H, \hat{H}]
\]

Da der Hamiltonoperator \(\hat{H}\) im Heisenberg-Bild zeitunabhängig ist, bleibt \(\hat{H}_H = \hat{H}\).